\subsection{The two character spaces at the next mixed degree}
\label{sec:kt-cubic-cohomology}

Put $\psi=\omega^e$, where $p>3$, $e\ge2$ is even, and
$(p-1)\nmid2e$.  Let $D<0$ be fundamental and prime to $p$,
$\eta=\chi_D$, and
$S=\{\infty,p\}\cup\{q:q\mid D\}$.  All cohomology below uses
this exact finite set.  The coefficient field $k$ is a finite extension
of $\mathbf F_p$; the trace pairing identifies the Cartier dual of
$k(\rho)$ with $k(\rho^{-1}\omega)$.

\begin{lemma}[KT-CUBIC-TWISTED-COHOMOLOGY]
If $B_{1,\eta\omega^e}$ is a $p$-adic unit, then
\[
 H^2(G_{\mathbf Q,S},k(\eta\psi^{-1}))=0,
 \qquad
 H^1(G_{\mathbf Q,S},k(\eta\psi^{-1}))
   \xrightarrow{\;\sim\;}
 H^1(\mathbf Q_p,k(\eta\psi^{-1})).                \tag{CC.1}
\]
Both $H^1$ spaces have dimension one.  Thus an arbitrary local
$\eta\psi^{-1}$-valued cocycle class at $p$ can be corrected by a
global cocycle in that character, uniquely at the level of classes.
\end{lemma}
\begin{proof}
It suffices to use $k=\mathbf F_p$ and then extend scalars.
Set $L=\mathbf Q(\sqrt D,\mu_p)$ and
$\Delta=\operatorname{Gal}(L/\mathbf Q)$.
For $M=\mathbf F_p(\eta\psi^{-1})$ the Cartier dual is
$M^D=\mathbf F_p(\eta\omega^{e+1})$.
The even mirror in KT-AUXILIARY-UNIT-SELMER-BRIDGE gives
\[
          \operatorname{Cl}(L)\{p\}^{\eta\omega^{e+1}}=0.
                                                               \tag{CC.2}
\]
A class in $\Sha^1_S(\mathbf Q,M^D)$ restricts to an everywhere
unramified homomorphism on $G_L$.  Since $|\Delta|=2(p-1)$ is
prime to $p$, restriction and class field theory embed this group
in $\operatorname{Hom}_\Delta(\operatorname{Cl}(L)/p,M^D)$.
Its target is zero by (CC.2); the character seen by equivariance
is $\eta\omega^{e+1}$ itself, not its inverse.
Poitou--Tate gives $\Sha^2_S(\mathbf Q,M)=0$.

Every local $H^2$ in $S$ vanishes.  At $p$, the dual inertia
character is $\omega^{e+1}$, which is nontrivial because $e+1$
is odd and $p-1$ is even.  At $q\mid D$ the dual has the ramified
quadratic inertia character $\eta$.  At infinity the higher
cohomology vanishes because $p$ is odd.  Local duality and the
Poitou--Tate exact sequence now prove the first assertion in
(CC.1).

The character $\eta\psi^{-1}$ is odd and nontrivial.  Global
Euler characteristic gives $h^1=1$ after the just-proved $h^2=0$.
At $p$ the coefficient character itself is nontrivial on inertia
because $e\not\equiv0\pmod{p-1}$; its local $H^0$ and $H^2$
are zero, so local Euler characteristic gives $h^1=1$.
At every $q\mid D$, both local $H^0$ and $H^2$ vanish and
$q\ne p$, so local $H^1$ is zero.
Consequently the kernel of localization at $p$ is contained in
the residual strict group killed by the same auxiliary-unit
bridge.  Localization is injective, and the equal dimensions
make it an isomorphism.
\end{proof}

\begin{lemma}[KT-CUBIC-RATIONAL-OBSTRUCTION-SPACE]
Let $F=\mathbf Q(\mu_p)$, $\Delta_0=\operatorname{Gal}(F/\mathbf Q)$,
$A_F=\operatorname{Cl}(F)\{p\}$, and $\delta=\omega^{1-e}$.
Let $C_D$ be the subgroup of $A_F$ generated by the $p$-primary
projections of the ideal classes of all primes of $F$ over $q\mid D$.
There is an exact sequence of finite $k$-vector spaces
\[
\begin{split}
0\longrightarrow
 \operatorname{Hom}_{\Delta_0}
   (A_F/(pA_F+C_D),k(\delta))^\vee
\longrightarrow H^2(G_{\mathbf Q,S},k(\psi))
\longrightarrow
 \bigoplus_{\substack{q\mid D\\q^{e-1}\equiv1\ (p)}}k
\longrightarrow0 .
\end{split}                                               \tag{CC.3}
\]
The left term is $\Sha^2_S(\mathbf Q,k(\psi))$.
Here $\vee$ denotes the linear dual, using the trace pairing
to express the finite-module duality over $k$.
\end{lemma}
\begin{proof}
The Cartier dual of $k(\psi)$ is the rational odd character
$k(\delta)$.  Restriction to $G_F$ identifies
$\Sha^1_S(\mathbf Q,k(\delta))$ with the displayed equivariant
Hom space before taking its dual.  To check equality, a strict
class becomes an everywhere unramified homomorphism on $G_F$,
and its vanishing on every decomposition group in $S$ kills the
classes of the primes over $D$.  The unique prime over $p$ is
principal in $F$, so it imposes no further class-group quotient.
Conversely a homomorphism from that quotient extends uniquely
as a cohomology class over $\mathbf Q$ by prime-to-$p$
inflation--restriction.  Its restrictions to the decomposition
groups in $S$ vanish by the same quotient and the principality
at $p$.  All local restriction maps used here are injective
under the extension of degree dividing $p-1$.

Poitou--Tate identifies the dual of that $\Sha^1$ with
$\Sha^2_S(\mathbf Q,k(\psi))$.  Since $1-e$ is odd, the
character $\delta$ is nontrivial both globally and on inertia
at $p$.  Thus its global $H^0$ and the local $H^2$ at $p$
vanish.  At $q\mid D$, $\delta$ is unramified and its Frobenius
value is $q^{1-e}$, so local duality gives a one-dimensional
$H^2$ exactly when $q^{e-1}\equiv1\pmod p$, and zero otherwise.
The real place contributes zero.  The last part of the
Poitou--Tate sequence, whose final term is the zero global
$H^0$ of the dual, proves (CC.3).
\end{proof}

The three auxiliary units at $0,e,2e$ annihilate the specified
$\eta$-twisted class components.  They supply (CC.1), but they
do not supply a vanishing assertion for the rational character
$\delta$ in (CC.3).  In particular, a raw global $H^2$ calculation
must also retain the displayed local terms.  A particular mixed
obstruction known to extend over every prescribed local functor
maps to zero in those local terms and belongs to the left term;
this observation does not prove that its remaining class is
zero or nonzero.

\paragraph{A finite check of the raw-local distinction.}
For $p=7,e=2,D=-43$, the defining Bernoulli sums give
\[
 B_{1,\eta}=-1,\qquad
 B_{1,\eta\omega^2}\equiv5\pmod7,\qquad
 B_{1,\eta\omega^4}\equiv3\pmod7.                    \tag{CC.4}
\]
Thus all three auxiliary units hold, but $43\equiv1\pmod7$
gives a nonzero local quotient in (CC.3).  In particular the
raw $H^2(G_{\mathbf Q,S},\mathbf F_7(\omega^2))$ is nonzero.
The separate standard-library certificate
\texttt{lean\_chain/cubic\_cohomology\_example.py} evaluates the
complete conductor sums modulo $7^2$ and independently checks
the two residues by exact rational Bernoulli polynomials of
indices $3,5$.  This example makes no claim that a specific
Massey obstruction or a strict Selmer group is nonzero.

\paragraph{Exact imported source scope.}
The unit-to-class and strict-cohomology bridge is the previously
whole-reviewed KT-AUXILIARY-UNIT-SELMER-BRIDGE.  The duality
inputs are Milne, \emph{Arithmetic Duality Theorems}, second
edition, Chapter I: the Cartier dual and $\Sha$ definitions
p.56, Theorem 4.10(a)--(b) pp.56--57 and its full exact sequence,
local Corollary 2.3 p.28, local Theorem 2.8 p.31, and global
Theorem 5.1 / Remark 5.2(c) pp.67--68.  Their complete statements
and exact finite-$S$ scope were independently read; the
foundational duality proofs remain imported.  These lemmas
describe the relevant cohomology spaces, and do not assert
that every abstract element is realized by the source's
higher ordinary deformation equations.
